%-*-latex-*-
%-*-latex-*-
\newcommand\COURSE{ciss350}
\newcommand\ASSESSMENT{t01}
\newcommand\ASSESSMENTTYPE{Test}
\newcommand\POINTS{\textwhite{xxx/xxx}}

\input{myquizpreamble}
\input{yliow}
\input{\COURSE}
\input{thispackages}
\input{thismacros}
%-*-latex-*-
\renewcommand\TITLE{\ASSESSMENTTYPE \ \ASSESSMENT}

\renewcommand\EMAIL{}
\renewcommand\AUTHOR{YOUR EMAIL}

\textwidth=6in
\begin{document}
\topmatterthree


%-*-latex-*-
\textsc{Instructions}
\begin{enumerate}
\li This is a closed-book, no-discussion, no-calculator, no-browsing-on-the-web
    no-compiler test.
\li Cheating is a serious academic offense. If caught you will 
    receive an immediate score of -100\%.
\li If a question asks for a program output and the program or
    code fragment contains
    an error, write \verb!ERROR! as output.
    When writing output, whitespace is significant.
\li If a question asks for the computation of a value and the program or
    code fragment contains
    an error, write \verb!ERROR! as value.
\li When you're asked to write a C++ statement, don't forget that it must
    end with a semicolon.
\li Unless otherwise stated, bubblesort, selectionsort,
    insertionsort, mergesort, and quicksort refers to the
    algorithms mention in class.
    Unless otherwise stated, sorting is always in ascending order.
\end{enumerate}

\vspace{1cm}


To enter your answers, look for \verb!answerbox! or \verb!answerlong!
in \verb!main.tex!.
If you have problems compiling, let me know.
When you are done:
\begin{enumerate}
  \li Submit using alex.
  \li Email me \verb!yliow@ccis.edu! with \verb!submit.tar.gz! attached and with
  subject line \verb!"ciss350 t01"!.
\end{enumerate}

\mbox{}\\

\begin{center}
  \textsc{Honor Statement}
\end{center}
I, \answerbox{[Ethan A. Graham]},
attest to the fact that the submitted work is my own and
is not the result of plagiarism.
Furthermore, I have not aided another student in the act of
plagiarism.



% ------------------------------------------------------------------------------
\begin{python}
from scoretable import *
\end{python}

%------------------------------------------------------------------------------
\newpage
%-*-latex-*-
\textsc{Some \LaTeX}

For $T(n) = (t_0 + t_1)n^2 + t_2 n + t_3$
write \verb!$T(n) = (t_0 + t_1)n^{2} + t_2n + t_3$!.

For $T(n) = O(n^{100})$
write \verb!$T(n) = O(n^{100})$!.

Let me know if you have questions about \LaTeX.

Note that $T(n) = n^2$ is not the same as $T(n) = O(n^2)$.



%------------------------------------------------------------------------------
\newpage
\nextq %[Derive big--O of runtime of a given algorithm]

The following computes the sum of squares from $1^2$ to $n^2$:
\begin{Verbatim}[frame=single]
s = 0
for i = 1, ..., n:
    term = i * i
    s = s + term
\end{Verbatim}
Here's the program with goto statements and timing for each statement:
\begin{Verbatim}[frame=single]
         i = 1              time t0
         s = 0              time t1
LOOP:    if i > n:          time t2
             goto ENDLOOP   time t3
         term = i * i       time t4
         s = s + term       time t5
         i = i + 1          time t6
         goto LOOP          time t7
ENDLOOP:
\end{Verbatim}
(a) Compute the time taken as a function of $n$.
The expression should contain constants $t_0$, ..., $t_7$.
\\
\ANSWER
\begin{answerlong}
$T(n) = (t_{2} + t_{4} + t_{5} + t_{6} + t_{7})n + (t_{0} + t_{1} + t_{3})$
\end{answerlong}

(b) Simplify the runtime function by giving names $A$, $B$, ...
to the constants of the function from (a).
\\
\ANSWER
\begin{answerlong}
$T(n) = An + B$
\end{answerlong}

(c) Fudge away the constants and write down the simplest $g(n)$ such 
that the time in (b) is a big-O of your $g(n)$.
Your $g(n)$ should be either $n^0$ (i.e., 1) or $n$ or $n^2$ or $n^3$ or ...
\\
\ANSWER
\begin{answerlong}
$T(n) = O(n)$
\end{answerlong}

% ------------------------------------------------------------------------------
\newpage
\nextq %[Recognize big--O of algorithm]

If the asymptotic runtime cannot be computed based on only what is taught in CISS350,
write ERROR in the answer.
Note that the algorithms are written in pseudocode (as usual).

(a)
What is the big--O of the runtime of the following algorithm:
\begin{Verbatim}[frame=single]
int f(int x[], int n)
{
    s = 0
    for i = 0, 1, 2, ..., n - 1:
        for j = 0, 1, 2, ..., i:
            for k = j, j + 1, ..., n - 1:
                s = s + x[i][j] + k
    return s
}
\end{Verbatim}
\ANSWER
\begin{answerlong}
$T(n) = O(n)^3$
\end{answerlong}


\mbox{}\\
(b)
What is the big--O of the runtime of the following algorithm:
\begin{Verbatim}[frame=single]
int f(int x[], int n)
{
    s = 0
    for i = 0, 1, 2, ..., n - 1:
        for j = n - 3, n - 2, n - 1, n:
            for k = i, i + 1, ..., n - 1:
                s = s + x[i][k] * j
    return s
}
\end{Verbatim}
\ANSWER
\begin{answerlong}
$T(n) = O(n)^2$
\end{answerlong}


\mbox{}\\
(c)
What is the big--O of the best runtime of the following algorithm:
\begin{Verbatim}[frame=single]
int f(int x[], int n)
{
    s = 0
    for i = 0, 1, 2, ..., n - 1:
        for j = n/2, n/2 + 1, ..., n:
            if x[0][0] > 0:
                for k = i, i + 1, ..., n - 1:
                    s = x[j][k] + x[i][k]
    return s
}
\end{Verbatim}
\ANSWER
\begin{answerlong}
$T(n) = O(n)$
\end{answerlong}


\mbox{}\\
(d)
What is the big--O runtime of the following algorithm:
\begin{Verbatim}[frame=single]
int f(int x[], int start, int end)
{
    n = end - start
    if n <= 1:
        return 42 + x[start]
    else:
        if x[0] == 42:
            return 24
        else:
            mid = start + n / 2
            if x[0] < 42:
                return f(x, start, mid)
            else:
                return f(x, mid, end) 
}
\end{Verbatim}
\ANSWER
\begin{answerlong}
$T(n) = O(n log n)$
\end{answerlong}


% ------------------------------------------------------------------------------
\newpage
\nextq %[Recall big-O runtime performance of standard algorithms]
For questions on runtime,
\underline{just state} the big-O of the runtime performance.

(a)
What is the runtime of insertion sort?
\\
\ANSWER
\begin{answerlong}
$T(n) = O(n)^2$
\end{answerlong}
  
(b)
What is the worst runtime of mergesort?
\\
\ANSWER
\begin{answerlong}
$T(n) = O(n log n)$
\end{answerlong}

(c)
What is the best runtime of mergesort?
\\
\ANSWER
\begin{answerlong}
$T(n) = O(n log n)$
\end{answerlong}
  
(d)
What is the average runtime of mergesort?
\\
\ANSWER
\begin{answerlong}
$T(n) = O(n log n)$
\end{answerlong}
  
(e)
What is the worst runtime of quicksort?
\\
\ANSWER
\begin{answerlong}
$T(n) = O(n)$
\end{answerlong}
  
(f)
What is the best runtime of quicksort?
\\
\ANSWER
\begin{answerlong}
$T(n) = O(n log n)$
\end{answerlong}
  
(g)
What is the average runtime of quicksort?
\\
\ANSWER
\begin{answerlong}
$T(n) = O(n log n)$
\end{answerlong}
  
(h)
What algorithm would you use to sort an array of 1000000000 doubles
with values $> 10$ and $< 1000000$?
(Assume you want the fastest of course.)
\\
\ANSWER
\begin{answerlong}
mergesort
\end{answerlong}
 
% ------------------------------------------------------------------------------
\newpage
\nextq %[Trace of basic sorting algorithms]
Perform insertion sort on the following array
\begin{console}
{5, 1, 2, 4, 3, 6}
\end{console}
showing the result of each pass (i.e., after each insert).
\\
\ANSWER
\begin{answercode}
{1, 5, 2, 4, 3, 6}
{1, 2, 5, 4, 3, 6}
{1, 2, 4, 5, 3, 6}
{1, 2, 4, 3, 5, 6}
{1, 2, 3, 4, 5, 6}
\end{answercode}
(You need not use all of the above arrays. Remove those not used.
Add more arrays if you need to.)

% ------------------------------------------------------------------------------
\newpage
\nextq
You are given the following singly linked node class:
\vspace{-3mm}
\begin{Verbatim}[frame=single,fontsize=\small]
class SLNode
{
public:
    SLNode(int key, SLNode * next=NULL)
    : key_(key), next_(next)
    {}
    
    int key_;
    SLNode * next_;
};
\end{Verbatim}
\vspace{-4mm}
Note that all members are public (to make life simple).
The following is an incomplete singly listed list class for integer
values,
\verb!SLList!.
Complete the following methods \verb!insert_head()! and \verb!delete_head()!.
Assume there is an \verb!UnderflowError! class that can be used for exception
throwing under a suitable condition.
\vspace{-3mm}
\begin{Verbatim}[frame=single,fontsize=\small]
class SLList
{
public:
    SLList()
      : phead_(NULL)
    {}
    void insert_head(int key);
    void delete_head();
    // ... other methods
private:
    SLNode * phead_;
};
\end{Verbatim}

(a)
\ANSWER
\begin{answercode}
void SLList::insert_head(int key)
{
   phead_ = new SLNode(*key, phead_);
}
\end{answercode}

(b)
\ANSWER
\begin{answercode}
void SLList::delete_head()
{
  SLNode temp = *phead_;
  phead_ = phead_ ->next;
  delete[] temp;
}
\end{answercode}

% ------------------------------------------------------------------------------
\newpage
\nextq %[Runtimes of singly linked list operations]
State the big-O of the runtime of the following operations
in a singly linked list ($n$ is the size of the list).

(a)
Insert head:
\\
\ANSWER
\begin{answerlong}
$T(n) = O(1)$
\end{answerlong}


(b)
Delete head:
\\
\ANSWER
\begin{answerlong}
$T(n) = O(1)$
\end{answerlong}
    
(c)
Insert tail:
\\
\ANSWER
\begin{answerlong}
$T(n) = O(n)$
\end{answerlong}

(d)
Delete tail:
\\
\ANSWER
\begin{answerlong}
$T(n) = O(n)?$
\end{answerlong}

(e)
Delete first node with a given key value:
\\
\ANSWER
\begin{answerlong}
$T(n) = O(n)$
\end{answerlong}
  
(f)
Compute the maximum of all the key values in the list:
\\
\ANSWER
\begin{answerlong}
$T(n) = O(n)$
\end{answerlong}
  
(g)
Swap the keys of the first two nodes (if there are at least two nodes):
\\
\ANSWER
\begin{answerlong}
$T(n) = O(1)$
\end{answerlong}

(h)
Insert a given key after the node of a given pointer:
\\
\ANSWER
\begin{answerlong}
$T(n) = O(n)$
\end{answerlong}

(i)
Insert a given key before the node of a given pointer:
\\
\ANSWER
\begin{answerlong}
$T(n) = O(n)$
\end{answerlong}

% ------------------------------------------------------------------------------
\newpage
\nextq %[Quicksort]
Assume the following functions for performing quicksort are
provided:
\vspace{-2mm}
\begin{Verbatim}[frame=single,fontsize=\small]
// Selects the index of a pivot in x[start .. end - 1]
int pivot_index(int x[], int start, int end);
 
// Performs a partition algorithm on x[start .. end - 1] with the pivot at
// x[i]. The return value is the index of the pivot at the end of the
// partition algorithm. 
int partition(int x[], int start, int end, int i);
\end{Verbatim}
\vspace{-4mm}
Complete the quicksort below.

\ANSWER
\begin{answercode}
void quicksort(int x[], int start, int end)
{ 
 int pivot = pivot_index(int x, start, end);
 
}
\end{answercode}

% ------------------------------------------------------------------------------
\newpage
\nextq %[Quicksort]
You are given the following array:
\begin{console}
{3, 4, 0, 1, 7, 6, 8, 9, 2}
\end{console}

(a)
Write down the array after performing the median-of-three
pivot selection and write down the pivot value.
\\
\ANSWER
\begin{answercode}
{2, 3, 0, 1, 4, 6, 8, 9, 7}
\end{answercode}
Pivot value is
\begin{answercode}
4
\end{answercode}


(b)
Write down the array after performing one partitioning.
(As mentioned in class, there are many possible partitioning
algorithms. You must use the method I that mentioned in class,
the method I called pLRT.)
\\
\ANSWER
\begin{answercode}
{0, 1, 2, 3, 4, 6, 8, 9, 7}
\end{answercode}

% ------------------------------------------------------------------------------
\newpage
\nextq
An array $x$ of size $n$ contains \verb!double!s
-2.0, -1.9, -1.8, ..., 0.1, 0, 0.1, ..., 2.0.
Note that $n$ is huge, approximately $1000000000$.
Explain how you would sort this array.
Note that you want the fastest possible algorithm.
Explain briefly how you would sort this array and provide an implementation below
(the name of the function is \verb!sort!).
In your explanation, include the runtime and how much extra memory you need.
The extra memory does not include the array and the size of the array.
But everything must be stated, even index variables used in loops. 
Complete the following sorting function.

\ANSWER

(a) \textsc{Explanation:}\vspace{-3mm}
\begin{answerlong}
... explanation ...
\end{answerlong}

(b) \textsc{Implementation:}\vspace{-2mm}
\begin{answercode}
void sort(double x[], int n)
{  
}  
\end{answercode}

% ------------------------------------------------------------------------------
\newpage
\nextq
Complete the following mergesort and merge functions.
Assume that the given temporary array is sufficiently large for
the merge function.

\ANSWER
\begin{answercode}
int t[1000000000]; // Temporary array that can be used by merge()

// Performs the merge operation on x[start0 .. end0 - 1] and
// x[start1 .. end1 - 1]. Here start1 = end0. You can assume the left and right
// subarray can fit into the above temporary array t.
void merge(int x[], int start0, int end0, int start1, int end1)
{
  
}

// Performs mergesort on x[start..end - 1]
void mergesort(int x[], int start, int end)
{
 if(start >= end)
   return;
 else
 {
   int mid = start + (start - end) / 2;
   mergesort(x, start, mid);
   mergesort(x, mid + 1, end);
   merge(x, start, mid, mid + 1, end);
 }
}
\end{answercode}

% ------------------------------------------------------------------------------
\newpage
\nextq % Quicksort
The following array has gone through one complete
partitioning step of the quicksort algorithm:
\begin{console}
{3, 2, 0, 4, 1, 5, 8, 6, 9}
\end{console}
The partitioning algorithm and the pivot selection algorithm
used are not known and might not be the ones you have seen before.
Write down all values which might be pivots?
\\
\ANSWER
\begin{answercode}
5, 9
\end{answercode}
(Edit the above. I'm not claiming there are exactly 3 possible pivots.)

% ------------------------------------------------------------------------------
\newpage
\nextq
Mergesort is performed on the following
array:
\begin{console}
{4, 2, 0, 3, 1, 5, 8, 6}
\end{console}

Of course we know that mergesort is recursive and calls itself.
When we \textit{return} from the \textit{first} mergesort function call,
the array is completed sorted.
What are the values of the array when we \textit{return} from the \textit{third}
mergesort function call.

\ANSWER
\begin{answercode}
{2, 4, 0, 3, 1, 5, 8, 6}
\end{answercode}

% ------------------------------------------------------------------------------
\newpage
\nextq
Suppose it is known that the array \verb!x! has even indexed values which are
ascending and odd indexed values which are ascending as well.
For instance in the array \verb!{5,0,8,1,9,2,10,7}!,
the values at even indexed positions are
\verb!5,8,9,10! which are in ascending order;
the values at odd indexed positions are
\verb!0,1,2,7! which are also in ascending order.
The array can have a huge size and the values need not be a small range.

How would you sort such an array?
(Of course you want the fastest algorithm that uses the least amount
of extra space, with speed being more important.)

Explain briefly (and clearly) your method.
Include your runtime and space complexity.
Next, implement the sorting function; the name of the function is \verb!sort!.
A temporary array \verb!t! (using \verb!std::vector!) is declared in the sorting
function.
If you do not need it, you should remove it.

\ANSWER

\textsc{Explanation}:
\begin{answerlong}

... explanation ...

The runtime is $O(?)$ and the space complexity is $O(?)$.
\end{answerlong}  

\textsc{Implementation}:
\begin{answercode} 
// n is the size of the array
void sort(int x[], int n)
{
    std::vector< int > t(100); // Here's a temporary array of size 100.
                               // Change 100 if necessary.
                               // If you do not need t, remove this statement.

} 
\end{answercode}

% ------------------------------------------------------------------------------
\newpage
\input{instructions.tex}
\end{document}
